% Môn: Vật lý
% Lớp: 12
% Chương: Chương I. Vật lí nhiệt
% Bài: L12C1 Bài 6. Nhiệt hóa hơi riêng · Bài toán nhiệt học tổng hợp yêu cầu tính năng lượng cho toàn bộ quá trình đun nóng khối chất lỏng đến nhiệt độ sôi rồi h
% ===== Câu 57 =====
% ID: L12C1B6-191-TN
% Mức: VD
\dangbt{Bài toán nhiệt học tổng hợp yêu cầu tính năng lượng cho toàn bộ quá trình đun nóng khối chất lỏng đến nhiệt độ sôi rồi h}
\begin{ex}
% ID: L12C1B6-191-TN
% Mức: VD
Tính nhiệt lượng cần cung cấp để đun $1,5\,\mathrm{kg}$ nước từ $20\,^\circ\mathrm{C}$ đến $100\,^\circ\mathrm{C}$ rồi hóa hơi hoàn toàn. Biết $c=4200\,\mathrm{J/(kg\cdot K)}$ và $L=2,3\cdot 10^6\,\mathrm{J/kg}$.
\choice
{$3450\,\mathrm{kJ}$}
{$504\,\mathrm{kJ}$}
{\True $3954\,\mathrm{kJ}$}
{$2954\,\mathrm{kJ}$}
\loigiai{
Nhiệt lượng đun nóng nước:
 $Q_1=mc\Delta t=1,5\cdot 4200\cdot 80=504000 \mathrm{J}=504 \mathrm{kJ}.$
 Nhiệt lượng hóa hơi:
 $Q_2=Lm=2,3\cdot 10^6\cdot 1,5=3450 \mathrm{kJ}.$
 Tổng nhiệt lượng:
 $Q=Q_1+Q_2=504+3450=3954 \mathrm{kJ}.$
}
\end{ex}

% ===== Câu 61 =====
% ID: L12C1B6-192-TN
% Mức: VD
\begin{ex}
% ID: L12C1B6-192-TN
% Mức: VD
Đun $0,8\,\mathrm{kg}$ rượu từ $25\,^\circ\mathrm{C}$ đến nhiệt độ sôi $78\,^\circ\mathrm{C}$ rồi làm rượu hóa hơi hoàn toàn. Biết $c=2400\,\mathrm{J/(kg\cdot K)}$ và $L=8,5\cdot 10^5\,\mathrm{J/kg}$. Nhiệt lượng cần cung cấp gần đúng là
\choice
{$680\,\mathrm{kJ}$}
{$102\,\mathrm{kJ}$}
{\True $782\,\mathrm{kJ}$}
{$952\,\mathrm{kJ}$}
\loigiai{
Nhiệt lượng đun nóng rượu:
 $Q_1=mc\Delta t=0,8\cdot 2400\cdot (78-25)=101760 \mathrm{J}\approx 102 \mathrm{kJ}.$
 Nhiệt lượng hóa hơi:
 $Q_2=Lm=8,5\cdot 10^5\cdot 0,8=680000 \mathrm{J}=680 \mathrm{kJ}.$
 Tổng nhiệt lượng:
 $Q=Q_1+Q_2\approx 102+680=782 \mathrm{kJ}.$
}
\end{ex}

% ===== Câu 65 =====
% ID: L12C1B6-193-TN
% Mức: VDC
\begin{ex}
% ID: L12C1B6-193-TN
% Mức: VDC
Một bình chứa $0,5\,\mathrm{kg}$ nước ở $30\,^\circ\mathrm{C}$. Người ta cung cấp cho nước nhiệt lượng $500\,\mathrm{kJ}$. Biết $c=4200\,\mathrm{J/(kg\cdot K)}$ và $L=2,3\cdot 10^6\,\mathrm{J/kg}$. Khối lượng nước hóa hơi gần đúng là
\choice
{$0,10\,\mathrm{kg}$}
{\True $0,154\,\mathrm{kg}$}
{$0,217\,\mathrm{kg}$}
{$0,50\,\mathrm{kg}$}
\loigiai{
Nhiệt lượng cần để đun nước từ $30\,^\circ\mathrm{C}$ đến $100\,^\circ\mathrm{C}$:
 $Q_1=mc\Delta t=0,5\cdot 4200\cdot 70=147000 \mathrm{J}=147 \mathrm{kJ}.$
 Nhiệt lượng còn lại dùng để hóa hơi:
 $Q_2=500-147=353 \mathrm{kJ}=353000 \mathrm{J}.$
 Khối lượng nước hóa hơi:
 $m'=\dfrac{Q_2}{L}=\dfrac{353000}{2,3\cdot 10^6}\approx 0,154 \mathrm{kg}.$
}
\end{ex}

% ===== Câu 69 =====
% ID: L12C1B6-194-TN
% Mức: VDC
\begin{ex}
% ID: L12C1B6-194-TN
% Mức: VDC
Tính nhiệt lượng cần cung cấp để biến hoàn toàn $2\,\mathrm{kg}$ nước đá ở $-10\,^\circ\mathrm{C}$ thành hơi nước ở $100\,^\circ\mathrm{C}$. Biết $c_{\text{đá}}=2100\,\mathrm{J/(kg\cdot K)}$, $\lambda=3,4\cdot 10^5\,\mathrm{J/kg}$, $c_{\text{nước}}=4200\,\mathrm{J/(kg\cdot K)}$, $L=2,3\cdot 10^6\,\mathrm{J/kg}$.
\choice
{$4600\,\mathrm{kJ}$}
{$5482\,\mathrm{kJ}$}
{\True $6162\,\mathrm{kJ}$}
{$722\,\mathrm{kJ}$}
\loigiai{
Quá trình gồm bốn giai đoạn:
 $Q_1=2\cdot 2100\cdot 10=42 \mathrm{kJ},Q_2=2\cdot 3,4\cdot 10^5=680 \mathrm{kJ},Q_3=2\cdot 4200\cdot 100=840 \mathrm{kJ},Q_4=2\cdot$ 2,3 $\cdot$ 10^6=4600 $\mathrm{kJ}$. $Tổng nhiệt lượng:$ Q=42+680+840+4600=6162 $\mathrm{kJ}$.
}
\end{ex}

% ===== Câu 73 =====
% ID: L12C1B6-195-TN
% Mức: VDC
\begin{ex}
% ID: L12C1B6-195-TN
% Mức: VDC
Một bếp điện đun $1,2\,\mathrm{kg}$ nước từ $25\,^\circ\mathrm{C}$ đến $100\,^\circ\mathrm{C}$ rồi hóa hơi hoàn toàn. Biết $c=4200\,\mathrm{J/(kg\cdot K)}$, $L=2,3\cdot 10^6\,\mathrm{J/kg}$ và hiệu suất của bếp là $80\%$. Điện năng bếp tiêu thụ là
\choice
{$3138\,\mathrm{kJ}$}
{$2760\,\mathrm{kJ}$}
{\True $3922,5\,\mathrm{kJ}$}
{$2510,4\,\mathrm{kJ}$}
\loigiai{
Nhiệt lượng có ích:
 $Q_{\text{ích}}=mc\Delta t+Lm=1,2\cdot 4200\cdot 75+2,3\cdot 10^6\cdot 1,2.$ Q_{\text{ích}}=378000+2760000=3138000 $\mathrm{J}=3138\mathrm{kJ}$. $Do hiệu suất$ H=80 $\$ %=0,8 $, điện năng bếp tiêu thụ là:$ A=\dfrac{Q_{\text{ích}}}{H}= $\dfrac{3138}{0,8}=3922,5\mathrm{kJ}$.
}
\end{ex}
